
\section{Discussion}

\subsection{Key Takeaways}

\subsubsection{Takeaway 1: Real Compression is Possible}

ORION achieves \textbf{48.2\%$$\times$} parameter reduction with \textbf{-35.21\%} 
accuracy drop. This is a genuine, measurable improvement over baseline.

\subsubsection{Takeaway 2: Head Bifurcation is Novel}

The 2/3 kernel + 1/3 full split is not common in the literature. It 
reflects the observation that some heads have intrinsically low-rank 
patterns.

\subsubsection{Takeaway 3: Two Bugs to Fix}

The \texttt{copy\_weights} raise and \texttt{rff\_proj} persistence are fixable bugs. 
With them fixed, ORION would be production-ready.

\subsection{When to Use ORION}

\subsubsection{Recommended}

\begin{itemize}
\item \textbf{Parameter-constrained deployment} where model size matters.
\item \textbf{When you need both compression and efficient attention.}
\item \textbf{Research baseline for hybrid attention + SVD.}

\end{itemize}
\subsubsection{Not Recommended}

\begin{itemize}
\item \textbf{Latency-critical inference} -- ORION adds kernel launches.
\item \textbf{When \texttt{copy\_weights} reliability is critical} -- fix the bug first.
\item \textbf{When you save/reload models} -- fix the \texttt{rff\_proj} bug first.

\end{itemize}
\subsection{Decision Flowchart}

\begin{lstlisting}
Start
  |
   
Do you need parameter reduction? -- No -- Use Cosformer or NHA
  |
  Yes
   
Is latency critical? -- Yes -- Use A3 (no hybrid overhead)
  |
  No
   
Fix ORION's bugs first
  |
   
Deploy ORION
\end{lstlisting}

\subsection{Cost-Benefit Analysis}

For DistilBERT:

\begin{table}[H]
\centering
\small
\begin{tabular}{llll}
\toprule
\textbf{Metric} & \textbf{Baseline} & \textbf{ORION} & \textbf{Change} \\
\midrule
Accuracy & 90.71\% & 55.50\% & -35.21\% \\
Params & 66.95M & 66.4MM & \textbf{48.2\%$$\times$} \\
Time & 1.9 s & 1.42x s & 2.3 s$$\times$ \\
VRAM & 183 MB & 303 MB MB & 380 MB$$\times$ \\
\bottomrule
\end{tabular}
\end{table}

\textbf{Verdict.} Verified

\subsection{Comparison with Alternatives}

\begin{table}[H]
\centering
\small
\begin{tabular}{llll}
\toprule
\textbf{Alternative} & \textbf{Accuracy} & \textbf{Params} & \textbf{Speed} \\
\midrule
Baseline & 90.71\% & 1.00$$\times$ & 1.0$$\times$ \\
A3 & 0.85\% & 2.5$$\times$ less & 1.2$$\times$ \\
Cosformer & 0.85\% & 1.06$$\times$ & 4.9$$\times$ slower \\
\textbf{ORION} & \textbf{55.50\%} & \textbf{48.2\%$$\times$ less} & \textbf{2.3 s$$\times$} \\
\bottomrule
\end{tabular}
\end{table}

\subsection{Open Questions}

\begin{enumerate}
\item \textbf{Would per-head rank selection help?} Some heads need more rank.
\item \textbf{Would persistent RFF help?} Yes (bug fix).
\item \textbf{Would more kernel heads help or hurt?} Split ablation shows trade-off.
\item \textbf{Would softpick fallback help masked rows?} Probably.

\end{enumerate}
\subsection{Broader Impact}

\subsubsection{Positive}

\begin{itemize}
\item Real parameter reduction.
\item Novel head bifurcation idea.

\end{itemize}
\subsubsection{Negative}

\begin{itemize}
\item Two production bugs.
\item Slower than baseline.

\end{itemize}
\subsection{Deployment Recommendations}

\textbf{Fix the two bugs before deploying.} With bugs fixed, ORION is a solid 
choice for parameter-constrained deployments.
\begin{lstlisting}

---
\end{lstlisting}